

\chapter*{ABSTRACT}
\phantomsection
\addcontentsline{toc}{chapter}{ABSTRACT}

Urban traffic congestion at signalized intersections in Kathmandu is exacerbated 
by the use of conventional fixed-time signal control, which fails to respond to 
fluctuating traffic demand. This study investigates the effectiveness of a VISVAP-based 
semi-actuated signal control system at the Babarmahal intersection. Field traffic data, 
including turning movements and geometric characteristics, were collected and analyzed 
using Highway Capacity Manual procedures to evaluate existing operational conditions. 
A detailed microsimulation model was developed in PTV VISSIM and calibrated to replicate 
real traffic behavior. The existing Fixed-Time Signal Control (FTSC) was modeled and validated 
based on observed volume and queue conditions. The FTSC was optimized using VISSIM built in 
optimization. Furthermore, it was manually optimized taking the signal phasing plans obtained 
from VISSIM as a base to achieve the best results. Subsequently, a semi-actuated control logic 
was used for off-peak hour which was designed in VISVAP using detector-based inputs, and non-linear 
traffic flow rates to optimize signaling plans. Multiple simulation scenarios were conducted under 
peak and off-peak traffic conditions to compare system performance. Key performance indicators 
including average delay, queue length, travel time, and level of service were evaluated. Results 
indicated that the manually optimized FTSC outperformed the other scenarios for peak hour and semi-actuated 
system significantly reduces unnecessary green time allocation and improves intersection 
efficiency for off-peak. A substantial decrease in average delay and queue length was 
observed, along with smoother traffic progression. The adaptive nature of the system allows 
better utilization of available green time under variable demand conditions. The findings demonstrate 
that semi-actuated signal control offers a practical and effective solution for improving urban 
intersection performance in developing cities. This study supports the implementation of demand-responsive 
traffic signal systems for sustainable traffic management in Kathmandu.

\vspace{1em}
\noindent\textit{\textbf{Keywords:} Traffic Signal Control, Semi-Actuated Control, VISVAP, PTV VISSIM, Traffic Simulation, Intersection Performance, Kathmandu.}